\section{Численные методы поиска экстремума функций}
\label{sec:numerical-optimization}

\subsection{Постановка задачи}

Задача безусловной оптимизации:
\[
    f(x)\to\min,
    \qquad x\in\mathbb R^n.
\]
Точка $x_*$ является локальным минимумом, если существует окрестность, в которой
$f(x_*)\leq f(x)$. Для дифференцируемой функции необходимое условие внутреннего
экстремума:
\[
    \nabla f(x_*)=0.
\]
Если $f\in C^2$ и гессиан $\nabla^2 f(x_*)$ положительно определён, то $x_*$ ---
строгий локальный минимум. Для выпуклой функции любой локальный минимум является
глобальным.

Численный метод строит последовательность
\[
    x_{k+1}=x_k+\alpha_k p_k,
\]
где $p_k$ --- направление поиска, $\alpha_k$ --- длина шага.

\subsection{Одномерная минимизация}

Для унимодальной функции на $[a,b]$ применяются методы, не использующие
производные.

\textbf{Метод дихотомии.} Интервал сужается сравнением значений в двух близких
точках около середины. Длина интервала уменьшается примерно вдвое за итерацию.

\textbf{Метод золотого сечения.} Внутренние точки выбираются так, чтобы одна из
них переиспользовалась на следующей итерации. Коэффициент сжатия:
\[
    \rho=\frac{\sqrt5-1}{2}\approx0{,}618.
\]
За итерацию требуется одно новое вычисление функции.

\textbf{Метод парабол.} По трём точкам строится квадратичный интерполянт, а его
вершина используется как новое приближение. Метод быстрее при гладкой функции,
но требует защитных условий. Практический алгоритм Брента сочетает параболическую
интерполяцию с золотым сечением.

Если доступна производная, минимум можно искать как корень $f'(x)=0$ методами
бисекции, секущих или Ньютона.

\subsection{Градиентный спуск}

Направление наискорейшего локального убывания:
\[
    p_k=-\nabla f(x_k).
\]
Итерация:
\[
    x_{k+1}=x_k-\alpha_k\nabla f(x_k).
\]
Шаг может быть постоянным, найденным одномерным поиском или выбранным по
условиям Армихо и Вольфа.

Если $f$ выпукла и её градиент липшицев с константой $L$, то выбор постоянного
шага $0<\alpha\le 1/L$ гарантирует убывание функции и сходимость значений
$f(x_k)$ к минимуму (при существовании минимайзера). Для $\mu$-сильно выпуклой
функции с подходящим постоянным шагом сходимость линейная, а скорость зависит от
отношения $L/\mu$, то есть от обусловленности.
На вытянутых линиях уровня градиентный спуск медленно ``зигзагирует''.

\subsection{Метод Ньютона}

Разложение второго порядка приводит к системе
\[
    \nabla^2 f(x_k)p_k=-\nabla f(x_k),
    \qquad
    x_{k+1}=x_k+p_k.
\]
Если гессиан достаточно гладок (например, локально липшицев) и в точке минимума
он невырожден, то при достаточно близком начальном приближении метод имеет
квадратичную сходимость. Недостатки:
вычисление и факторизация гессиана, отсутствие глобальной сходимости и возможная
неположительная определённость вдали от минимума.

Глобализацию обеспечивают демпфирование
\[
    x_{k+1}=x_k+\alpha_kp_k
\]
или доверительная область, где квадратичная модель минимизируется при
$\|p\|\leq\Delta_k$.

\subsection{Квазиньютоновские методы}

Они строят приближение гессиана или обратного гессиана по изменениям
\[
    s_k=x_{k+1}-x_k,
    \qquad
    y_k=\nabla f(x_{k+1})-\nabla f(x_k).
\]
Наиболее распространён метод BFGS. Обновление обратного гессиана:
\[
H_{k+1}=
(I-\rho_ks_ky_k^T)H_k(I-\rho_ky_ks_k^T)
+\rho_ks_ks_k^T,
\qquad
\rho_k=\frac1{y_k^Ts_k}.
\]
При $y_k^Ts_k>0$ положительная определённость сохраняется. L-BFGS хранит только
несколько последних пар $(s_k,y_k)$ и подходит для задач большой размерности.

\subsection{Метод сопряжённых градиентов}

Для квадратичной функции
\[
    f(x)=\frac12x^TAx-b^Tx,
    \qquad A=A^T>0,
\]
метод строит $A$-сопряжённые направления и в точной арифметике достигает решения
не более чем за $n$ итераций. Для общего нелинейного функционала существуют
нелинейные варианты Флетчера--Ривса и Полака--Рибьера.

\subsection{Условная оптимизация}

Пусть
\[
    g_i(x)\leq0,
    \qquad
    h_j(x)=0.
\]
Функция Лагранжа:
\[
    \mathcal L(x,\lambda,\mu)
    =f(x)+\sum_i\lambda_i g_i(x)+\sum_j\mu_j h_j(x).
\]
При выполнении условий регулярности необходимы условия Каруша--Куна--Таккера:
\[
\begin{aligned}
&\nabla_x\mathcal L(x_*,\lambda_*,\mu_*)=0,\\
&g_i(x_*)\leq0,\qquad h_j(x_*)=0,\\
&\lambda_i\geq0,\\
&\lambda_i g_i(x_*)=0.
\end{aligned}
\]
В выпуклой задаче, когда $f$ и $g_i$ выпуклы, равенства $h_j$ аффинны и
выполнено условие Слейтера, условия KKT не только необходимы, но и достаточны
для глобального минимума.

Численные подходы:
\begin{itemize}
    \item проекция на допустимое множество;
    \item штрафные и барьерные функции;
    \item метод множителей Лагранжа;
    \item последовательное квадратичное программирование;
    \item методы внутренней точки.
\end{itemize}

\subsection{Критерии остановки}

Обычно контролируют несколько величин:
\[
    \|\nabla f(x_k)\|\leq\varepsilon_g,
    \qquad
    \|x_{k+1}-x_k\|\leq\varepsilon_x(1+\|x_k\|),
\]
\[
    |f(x_{k+1})-f(x_k)|\leq\varepsilon_f(1+|f(x_k)|).
\]
Для условной задачи дополнительно проверяются невязки ограничений и условия KKT.
Остановка только по малому изменению функции может быть ложной на плоском участке.

\subsection{Условия первого и второго порядка}

Пусть $f\in C^2$ и $x^*$ --- внутренняя точка локального минимума. Тогда необходимо
\[
    \nabla f(x^*)=0,
    \qquad
    \nabla^2f(x^*)\succeq0.
\]
Если
\[
    \nabla f(x^*)=0,
    \qquad
    \nabla^2f(x^*)\succ0,
\]
то $x^*$ является строгим локальным минимумом. Для выпуклой дифференцируемой функции
условие $\nabla f(x^*)=0$ уже достаточно для глобального минимума.

Для $m$-сильно выпуклой функции с $L$-липшицевым градиентом
\[
    mI\preceq \nabla^2f(x)\preceq LI
\]
в квадратичном смысле. Число $\kappa=L/m$ играет роль числа обусловленности задачи
оптимизации и определяет скорость работы методов первого порядка.

\subsection{Скорость градиентного метода}

При фиксированном шаге $\alpha=1/L$ для $m$-сильно выпуклой $L$-гладкой функции
градиентный метод
\[
    x_{k+1}=x_k-\frac1L\nabla f(x_k)
\]
сходится линейно. В одной из стандартных форм оценка записывается как
\[
    f(x_k)-f(x^*)
    \le
    \left(1-\frac{m}{L}\right)^k
    \bigl(f(x_0)-f(x^*)\bigr).
\]
Чем больше $L/m$, тем медленнее сходимость. Для квадратичной задачи это полностью
согласуется с ролью спектрального числа обусловленности матрицы Гессе.

\subsection{Линейный поиск: Армихо и Вульф}

Вместо заранее фиксированного шага выбирают направление $p_k$ и длину шага $\alpha_k$.
Условие Армихо требует достаточного уменьшения:
\[
 f(x_k+\alpha p_k)
 \le
 f(x_k)+c_1\alpha\nabla f(x_k)^Tp_k,
 \qquad 0<c_1<1.
\]
Условия Вульфа дополняют его требованием на производную вдоль направления. Такой
линейный поиск предотвращает слишком большие шаги и одновременно не даёт выбирать
бесконечно малые шаги без необходимости.

Для метода Ньютона с line search далеко от решения шаг можно уменьшать, а в окрестности
невырожденного минимума обычно принимается полный шаг $\alpha=1$, после чего проявляется
квадратичная сходимость.

\subsection{Квазиньютоновское условие и BFGS}

Квазиньютоновские методы строят приближение $B_k$ к гессиану или $H_k$ к его обратной
матрице, используя пары
\[
    s_k=x_{k+1}-x_k,
    \qquad
    y_k=\nabla f(x_{k+1})-\nabla f(x_k).
\]
Основное секущее условие:
\[
    B_{k+1}s_k=y_k.
\]
Формула BFGS для приближения гессиана
\[
 B_{k+1}
 =B_k-\frac{B_ks_ks_k^TB_k}{s_k^TB_ks_k}
 +\frac{y_ky_k^T}{y_k^Ts_k}
\]
сохраняет симметрию и при $y_k^Ts_k>0$ --- положительную определённость. Поэтому BFGS
с хорошим линейным поиском является одним из основных универсальных методов гладкой
безусловной оптимизации.

\subsection{Штрафы, барьеры и активные ограничения}

Для ограничений $g_i(x)\le0$ внешний штраф может иметь вид
\[
    F_\rho(x)=f(x)+\rho\sum_i\max\{0,g_i(x)\}^2,
    \qquad \rho\to\infty.
\]
Барьерный подход, напротив, остаётся внутри допустимой области, например
\[
    F_\mu(x)=f(x)-\mu\sum_i\ln[-g_i(x)],
    \qquad \mu\downarrow0.
\]
Современные interior-point методы следуют центральной траектории, определяемой
возмущёнными условиями KKT. Методы активного множества пытаются определить, какие
ограничения в решении выполняются как равенства, и затем решают соответствующую
локальную задачу с равенствами.

\subsection{Что важно сказать на экзамене}

Нужно разделить одномерные методы, методы первого и второго порядка и условную
оптимизацию. Градиентный спуск дёшев, но чувствителен к обусловленности; Ньютон
быстр вблизи решения, но дорог; BFGS даёт практический компромисс. Для ограничений
основным теоретическим инструментом являются условия KKT.

\newpage